\section{Contact and reciprocal selection of opposing regions}
\label{sec:area_law_reciprocal_opponents}

The unique opposing region determines a selection for each actual
elementary side. Contact with an actual fine cell characterizes this
selection. When the lower index is sufficiently large, the corresponding
elementary segments select each other; see
\cite[Section~11]{OpenAI2026AreaLaw}.

% Source: 10-geometry.tex, prop:two-families, lines 299--316 for
% choice/contact and lines 299--323 for reciprocal selection,
% revision adc7f1241b42e322a6451854ab7e4b4c146bf78a.
% Scoped statements only; all three released declarations independently reviewed.

\begin{definition}[The selected opposing region]
    \label{def:area_law_selected_elementary_opponent}
    \lean{TNLean.PEPS.AreaLaw.Geometry.elementarySideOpponent}
    \leanok
    \uses{thm:area_law_unique_elementary_side_opponent,
        def:area_law_actual_side_mask,def:area_law_fine_layer_indices}
    Fix a common origin $o\in\mathbb R^2$, finite endpoint set $Z$,
    integer width $C_0\ge2$, and natural indices $k_0\le k$ with
    $k\ge50\,000\,000$. For an actual fine-cell index
    $z\in F_{k,\ell_k}$ and a slot $i$ of its actual midpoint
    subdivision, write $e_i=[a_i,b_i]$.
    Denote by $P_{k,z}(i)$ the unique opposing identifier whose closed
    region contains $e_i$: the dummy symbol $\star$ represents
    $\overline{N_{k_0}}$, while an indexed cell $(h,w)$ represents
    $\overline{Q_{\ell_h}(w)}$ with $h\ge k_0$,
    $w\in F_{h,\ell_h}$ and $(h,w)\ne(k,z)$.
    The dummy neighborhood is one identifier even if several coarse
    cells contribute to it. No late bound on $k_0$ is required.
\end{definition}

\begin{theorem}[Selection is equivalent to contact with an actual fine cell]
    \label{thm:area_law_selected_opponent_contact}
    \lean{TNLean.PEPS.AreaLaw.Geometry.elementarySideOpponent_eq_some_iff_contact}
    \leanok
    \uses{def:area_law_selected_elementary_opponent,
        def:area_law_fine_layer_indices}
    Retain the assumptions of the selected-opponent definition.
    Let $w\in F_{h,\ell_h}$ be an actual candidate fine-cell index
    at a layer $h\ge k_0$. Then
    \begin{align}
        P_{k,z}(i)=(h,w)
        \quad\Longleftrightarrow\quad
        (k,z)\ne(h,w)\ \text{and}\
        e_i\cap\overline{Q_{\ell_h}(w)}
        \text{ contains two distinct points}.
    \end{align}
    The candidate layer $h$ has no additional late bound.
\end{theorem}
\begin{proof}\leanok
    \uses{thm:area_law_unique_elementary_side_opponent,
        thm:area_law_elementary_side_boundary_geometry,
        thm:area_law_actual_elementary_side_matching,
        thm:area_law_cell_fan_mark_vertices,
        thm:area_law_belt_marks_closed_cell,thm:area_law_dyadic_closures}
    If the candidate is selected, its indexed cell differs from the
    reference and its closure contains the whole segment. The two
    elementary endpoints are distinct, so both belong to the contact.

    Conversely, distinct indexed cells with two-point contact have
    matching actual elementary segments with reversed endpoints.
    The candidate segment lies in the closed candidate square:
    its endpoints belong to that square and the square is convex.
    Thus the whole reference segment lies in the same closed square.
    Uniqueness of the opposing identifier makes the candidate the
    selected opponent.
\end{proof}

\begin{theorem}[Reciprocal selection along matching elementary sides]
    \label{thm:area_law_reciprocal_elementary_opponents}
    \lean{TNLean.PEPS.AreaLaw.Geometry.elementarySideOpponent_reciprocal_iff}
    \leanok
    \uses{def:area_law_selected_elementary_opponent,
        def:area_law_actual_side_mask,def:area_law_fine_layer_indices}
    Fix a common origin and finite endpoint set, integer width
    $C_0\ge2$, and natural indices $k,h\ge k_0\ge50\,000\,000$.
    Let $z\in F_{k,\ell_k}$ and $w\in F_{h,\ell_h}$ be actual
    fine-cell indices. For every actual elementary slot $i$ of
    the reference cell, the equality $P_{k,z}(i)=(h,w)$ holds
    if and only if $(k,z)\ne(h,w)$ and there exists an actual
    elementary slot $j$ of the candidate cell such that
    \begin{align}
        s(j)       & =s(i)+2\pmod4,       \\
        a_i        & =b_j,\qquad b_i=a_j, \\
        P_{h,w}(j) & =(k,z).
    \end{align}
    Here $s(i)$ and $s(j)$ are the counterclockwise side indices,
    and both subdivisions use the same lower index $k_0$.
    No matching or return selection is supplied as a premise.
\end{theorem}
\begin{proof}\leanok
    \uses{thm:area_law_selected_opponent_contact,
        thm:area_law_actual_elementary_side_matching,
        thm:area_law_elementary_side_boundary_geometry,
        thm:area_law_cell_fan_mark_vertices,
        thm:area_law_belt_marks_closed_cell,thm:area_law_dyadic_closures}
    The lower bound on $k_0$ gives the required late bounds on
    both $k$ and $h$. If the candidate is selected, the contact
    characterization gives distinctness and two-point contact.
    Actual matching supplies the slot $j$, opposite side index
    and reversed endpoints. These endpoints belong to the
    closed reference square, which is convex; they are distinct.
    Applying the contact characterization with the cells exchanged
    shows that $j$ selects the reference cell in return.

    Conversely, the reversed endpoints make the two elementary
    segments equal. The candidate segment is nondegenerate and
    lies in its closed square. Hence the reference segment has
    two-point contact with that square. Distinctness and the
    contact characterization give $P_{k,z}(i)=(h,w)$.
\end{proof}

These assertions establish reciprocal selection of actual fine-cell
opponents. They do not assign primary identifiers or colors. Consistent
global labels, the isolated-star construction and the complete two-family
partition remain further obligations.
